\section{Problem Formulation}
\subsection{Optional-source molecular language modeling}
Let $\mathcal G$ be the space of molecular graphs and let a property program
be a variable-length list
\begin{equation}
  P=\{(j_\ell,g_\ell)\}_{\ell=1}^{m},
  \label{eq:property-program-unified}
\end{equation}
where $j_\ell$ names a property and $g_\ell$ is either a numerical target or a
directional goal. A request is $x=(P,S)$, where $S=\texttt{<EMPTY>}$ denotes de
novo construction and otherwise contains a canonical source SMILES for
editing. Both regimes are modeled by the same conditional distribution
\begin{equation}
  y=(z,s)\sim p_\theta(y\mid P,S),
  \label{eq:single-policy}
\end{equation}
where $z\in\{\texttt{BUILD},\texttt{MODIFY}\}$ and $s$ is a complete SMILES
string. The source field is an observed part of the request, not the output of
a learned router.

We call a system \emph{model-level unified} when the two regimes share the
same base weights, trainable adapter, tokenizer, prompt template, and output
schema. This definition excludes a router, a mode-specific adapter, and a
mode-specific output head. It does not require the external evaluator or
candidate-selection rule to be identical across benchmarks. The latter
distinction matters because editing includes source fidelity whereas de novo
construction does not.

\subsection{External verification and budgeted success}
Let $B$ denote a benchmark contract and
$V_B(P,G;S)\in\{0,1\}$ its strict verifier. It combines chemical validity and
joint property satisfaction; for editing it additionally requires source
Tanimoto similarity above the declared threshold. For an ordered prefix of
$K$ raw candidates, the editing endpoint is
\begin{equation}
  \operatorname{AnyAcc}^{@K}_{B}
  =\frac{1}{n}\sum_{i=1}^{n}
  \mathbf 1\!\left[\max_{1\le r\le K}
  V_B(P_i,G_i^{(r)};S_i)=1\right].
  \label{eq:any-accuracy}
\end{equation}
We use the more explicit notation
$\operatorname{AnyAcc}^{@K}_{\rm all}(.65)$ for the strict
MolEdit-Instruct endpoint. This is condition-level success within a budget and
must not be compared directly with candidate-level accuracy.

The matched SketchMol protocol instead draws 40 candidates and applies its
declared property-aware finalizer. We denote the resulting endpoint
$\operatorname{Best}_{B}^{@40}$. Targets and output properties are unavailable
during generation and enter only this external post-generation selection. We
therefore never pool the editing and de novo values into one numerical score.

\section{MolProgram}
\subsection{One serialization and one checkpoint}
MolProgram uses Qwen2.5-VL-7B-Instruct with a single LoRA adapter. The system
message instructs one molecular causal language model to infer the requested
operation only from the source field. The user message is deterministic JSON,
\begin{equation}
  \texttt{\{"conditions":[...],"source":S\}},
  \label{eq:input-schema}
\end{equation}
and the response must match exactly
\begin{equation}
  \texttt{\{"plan":"BUILD|MODIFY","smiles":"CANONICAL\_SMILES"\}}.
  \label{eq:output-schema}
\end{equation}
Only an allowlist of eight controls (MW, LogP, QED, TPSA, HBD, HBA, RB, and
SA) is serialized. Target molecules, target scaffolds, and oracle outputs are
excluded from the prompt. De novo and editing examples use the same chat
template and causal-LM output vocabulary; the plan token is supervised but
does not select another network.

The supervised checkpoint continues a frozen mixed de novo/editing LoRA. Its
training set contains 160 de novo and 160 editing examples, balanced at the
level of chosen-completion likelihood. Training uses one epoch at learning
rate $2\times10^{-5}$. The 320 unique examples expand to 800 logical
chosen--negative pairs, but each unique chosen completion contributes total
weight one regardless of how many negative types it has.

\subsection{Validity-aware multi-negative contrastive SFT}
For example $i$, let $\ell_i^+$ be the mean token negative log-likelihood of
the chosen completion and let $\ell_{in}^-$ be that of negative type $n$. Up
to a constant minibatch scaling, the implemented objective is
\begin{equation}
  \mathcal L_i=\ell_i^+
  +\sum_{n\in\mathcal N_i}\lambda_n
  \left[m_n+\ell_i^+-\ell_{in}^-\right]_+ .
  \label{eq:multinegative-objective}
\end{equation}
Each negative is forwarded separately. The hinge term is active when the
chosen completion is not preferred by at least margin $m_n$.

We use three train-only negative constructions. First, every example receives
an invalid corruption formed by appending one unmatched parenthesis to the
chosen SMILES ($m=0.15$, $\lambda=0.20$). Second, every example receives a
valid same-mode target from a different condition ($m=0.10$,
$\lambda=0.10$). Third, editing examples receive the unchanged source as a
weak copy negative ($m=0.10$, $\lambda=0.12$). The invalid strings retain the
strict JSON wrapper but are RDKit-invalid; mismatch donors never share the
same target and never equal the edit source. All negatives are constructed
from training rows only. This is multi-negative contrastive SFT with
ORPO-like margins, not ORPO and not reinforcement learning.

\subsection{Inference and benchmark boundary}
At inference, the model receives $(P,S)$ and emits an ordered stream from the
same checkpoint. No target molecule, target property value, static molecule
pool, or learned task router is available during decoding. For editing, we
report prefix Any@$K$ at $K\in\{1,4,8\}$ using the benchmark-native property
direction and source-similarity verifier. For de novo construction, we report
both raw-prefix diagnostics and the matched property-aware best-of-40
finalizer. Keeping the finalizer outside the model is why the present system
is model-level unified rather than fully end-to-end unified.

\subsection{Verifier-aligned group-relative post-training}
The supervised checkpoint exhibits a large gap between raw candidate success
and best-of-many success. We use this gap as evidence that useful molecules are
present in the policy support but underweighted, and continue the same adapter
with a critic-free group-relative policy update. We select 128 training
requests, balanced between construction and editing and across observed
property counts. For each request $x_i$, the current policy samples a group of
$K=4$ raw completions. Let $R_{ir}$ denote the terminal program reward for
completion $r$. We form within-request advantages
\begin{equation}
  A_{ir}=\frac{R_{ir}-\overline R_i}
  {\operatorname{Std}(R_{i1},\ldots,R_{iK})+\epsilon}
  \quad\text{with}\quad A_{ir}\in[-3,3],
  \label{eq:group-advantage}
\end{equation}
and minimize
\begin{equation}
  \mathcal L_{\rm post}
  =-\frac{1}{K}\sum_{r=1}^{K}A_{ir}
  \log p_\theta(y_{ir}\mid x_i)
  +\lambda_{\rm SFT}\mathcal L_{\rm SFT},
  \qquad \lambda_{\rm SFT}=0.20.
  \label{eq:group-policy-objective}
\end{equation}
The sampled-completion term uses mean completion log probability, and the SFT
anchor uses the training chosen completion. We perform one epoch at learning
rate $10^{-6}$. The update uses neither a learned critic nor PPO probability
ratios, clipping, or a reference-model KL penalty. It is therefore a
GRPO-style group-relative policy update, not an exact reproduction of the
clipped GRPO objective in DeepSeekMath \citep{shao2024deepseekmath}.

The reward exposes the same property-program structure in both regimes. An
invalid or schema-violating response receives $-1$. A valid response is scored
by canonicality, the fraction and mean degree of satisfied clauses, a smooth
bottleneck term derived from the least-satisfied clause, and an all-clause
bonus. Editing additionally receives continuous Morgan-fingerprint Tanimoto
reward, a bonus at similarity $\geq0.65$, a strict property-and-similarity
bonus, and a small exact-copy penalty. Target SMILES are never read by the
reward; they are used only by the SFT anchor. The post-training split contains
no benchmark rows, and inference retains the same one-checkpoint interface and
ordered raw candidate stream.
